\documentclass[11pt,a4paper]{article}
\usepackage{fontspec}
\setmainfont{Latin Modern Roman}
\usepackage[margin=23mm]{geometry}
\usepackage{amsmath,amssymb,amsthm,mathtools,mathrsfs}
\usepackage{longtable,booktabs}
\usepackage[unicode,hidelinks]{hyperref}
\providecommand{\C}{\mathbb C}
\providecommand{\R}{\mathbb R}
\newtheorem{theorem}{Theorem}[section]
\newtheorem{lemma}[theorem]{Lemma}
\newtheorem{proposition}[theorem]{Proposition}
\newtheorem{corollary}[theorem]{Corollary}
\newtheorem{definition}[theorem]{Definition}
\newtheorem{remark}[theorem]{Remark}
\allowdisplaybreaks
\emergencystretch=3em
\makeatletter
\renewcommand\@pnumwidth{2.5em}
\renewcommand\@tocrmarg{3.5em}
\makeatother
\title{The Original Quadric Orbit, Invariant Lifting and Exact Rank}
\author{Split-Zero research programme}
\date{14 September 2026}
\begin{document}
\maketitle
\paragraph{Portable source edition.}
Every TeX input used by this entry is expanded below. Historical file
and directory names in the preserved text are source locators; the
exact origin paths and hashes are recorded in \texttt{11\_SOURCE\_MAP.json}.
The other numbered proof files in this flat handoff supply the
accompanying complete calculations.\par

% Derived from the literal original period quadric and cyclic projector.
\section{The original quadric orbit and an explicit invariant lifting map}
\label{sec:OQC}

This calculation uses the original simple-quartet construction, with
$m=1$, $k=4l+1$, $l\geq1$, $q=(k+1)^2$, $0<\delta<1/2$, and
$\gamma>2$. The period parameter, its branch, the polynomial contours,
the complete amplitude unit, and the observation are those of OPG and
OCS. The degree-eight polynomial below is the product of four actual
orbit equations. Its degree and all of its coefficients are derived
from that orbit; no source order is selected to match an estimate.

\subsection{The actual coefficient ring and the original observation}
Order the original roots as $(++),(+-),(-+),(--)$, with
$\alpha_{\epsilon,\eta}=1/2+\epsilon\delta+i\eta\gamma$.
Let $\mathcal E_{\rm prim}$ have columns the four original primary
idempotents in $E_h=\mathbb C[z_1]/(h)$. For the original matrices
$\Pi,U,F$ in OCS2, OCS5, OCS27 put
\[
 \mathsf H=F^{-1}\Pi U\mathcal E_{\rm prim},\qquad
 t_{\epsilon,\eta}(w)=e^{(\eta\gamma-i\epsilon\delta)w},\qquad
 x(w)=\mathsf Ht(w),\qquad
 Q(x)=Q_0(\mathsf H^{-1}x),\quad
 Q_0(t)=t_{++}t_{--}-t_{+-}t_{-+}.
 \tag{OQC1}
\]
Every matrix in $\mathsf H$ is the specified original isomorphism;
in particular its nonzero amplitude factors are the actual
$v_h(\alpha)$, and its period entries are the original convergent
polynomial integrals. Thus $\det\mathsf H\ne0$ by the included
period-isomorphism and primary-basis proofs. Define
\[
 \mathscr P_d=\mathbb C[x_1,x_2,x_3,x_4]_d,\qquad
 R_d=\mathscr P_d/(Q\mathscr P_{d-2}),\qquad
 \mathscr P_d=R_d=0\quad(d<0).
 \tag{OQC2}
\]
The matrix of $Q_0(t)=t^TM_0t$ has entries
$(M_0)_{14}=(M_0)_{41}=1/2$ and
$(M_0)_{23}=(M_0)_{32}=-1/2$, all others zero. Hence
\[
 M=\mathsf H^{-T}M_0\mathsf H^{-1},\qquad
 \det M=\frac{1}{16(\det\mathsf H)^2}\ne0.
 \tag{OQC3}
\]
In particular $Q$ is irreducible. Indeed, a reducible homogeneous
quadratic over $\mathbb C$ is a product of two linear forms, whose
symmetric coefficient matrix is a sum of two rank-one matrices and
has rank at most two. This contradicts OQC3. The polynomial ring
over a field is a unique factorization domain, by the repeated
one-variable Gauss lemma, so its irreducible element $Q$ is prime.
Consequently $R=\mathbb C[x_1,x_2,x_3,x_4]/(Q)$ is an integral domain.

For clarity, the complete homogeneous restriction calculation has
the precise kernel $Q\mathscr P_{d-2}$. In the $t$ coordinates a
degree-$d$ monomial restricts to
\[
 \exp\bigl(((2b-d)\gamma-i(2a-d)\delta)w\bigr),
       \qquad 0\leq a,b\leq d.
 \tag{OQC4}
\]
Every pair occurs: the number of $(++)$ factors can range from
$\max(0,a+b-d)$ to $\min(a,b)$. The exponents are distinct by
their real and imaginary parts. Their exponentials are linearly
independent, since differentiating orders zero through $(d+1)^2-1$
gives a Vandermonde matrix with nonzero determinant. Two monomials
giving the same pair have exponent vectors differing by a multiple
of $(1,-1,-1,1)$. After their common monomial is removed, their
difference is divisible by $t_{++}t_{--}-t_{+-}t_{-+}$ through
$A^r-B^r=(A-B)\sum_{j=0}^{r-1}A^{r-1-j}B^j$. A vanishing
polynomial has coefficient sum zero within each pair by the
independence just proved, and is therefore a sum of these differences.
Conversely $Q_0(t(w))=0$. Transport by $\mathsf H$ proves the
stated kernel and
\[
                  \dim R_d=(d+1)^2\quad(d\geq0).
 \tag{OQC5}
\]

Write $\zeta=e^{2\pi i/5}$ and
$D=\operatorname{diag}(\zeta,\zeta^2,\zeta^3,\zeta^4)$.
The literal Fourier calculation gives $CF=FD$. On polynomials
define the original induced action and its average by
\[
 (g f)(x)=f(D^{-1}x),\qquad g^5=1,\qquad
 \mathsf R f=\frac15\sum_{a=0}^4g^af,
 \qquad \mathscr I_d=\mathscr P_d^{\langle g\rangle}.
 \tag{OQC6}
\]
Thus $\mathscr I_d$ consists of the weight-zero monomials, since
$gx^\nu=\zeta^{-\sum j\nu_j}x^\nu$. This is the dual of the
original average of $(C^{-a})^{\otimes d}$; both the inverse powers
and the factor $1/5$ are retained.

For the original algebra $E_d$ of the $d$-fold sum, define
$\Theta_d:R_d\xrightarrow{\sim}E_d^\vee$ by
\[
 \Theta_d([f])(e^{wY_d}[1])=f(x(w)),\qquad
       Y_d=(A_d-dI/2)/i.
 \tag{OQC7}
\]
Here $\vee$ denotes the complex-linear dual. The primary expansion
of the vector on the left gives all distinct exponentials OQC4;
therefore existence and uniqueness of this covector and bijectivity
follow from the preceding complete restriction calculation. The
original observation $\Lambda_d:E_d\to B_d$ obeys
\[
 J_d:=\operatorname{im}(\mathscr I_d\to R_d),\qquad
 \Theta_d(J_d)=\operatorname{im}\Lambda_d^\vee,
 \qquad \dim J_d=\operatorname{rank}\Lambda_d,
 \qquad R_d/J_d\xrightarrow{\sim}(\ker\Lambda_d)^\vee.
 \tag{OQC8}
\]
Indeed the coordinate functions of the projected pure tensor in
OCS9 are precisely $x(w)^\nu$ of weight zero; their pulled-back
covectors span the dual of its actual image. A covector vanishing
on $\ker\Lambda_d$ factors through that image, which proves the
last isomorphism as well. OQC8 specifies the exact original image;
it has not been replaced by the whole invariant tensor space.

\subsection{The complete orbit decision from the actual coefficients}
Put $Q_a=g^aQ$ for $0\leq a<5$. The orbit of the line $\mathbb C Q$
has size one or five because five is prime. Moreover a size-one
orbit has $gQ=Q$, with no other possible scalar. To prove this, write
$gQ=\lambda Q$. Then $\lambda^5=1$, while OQC3 and the determinant
of $D$, which is $\zeta^{10}=1$, imply
\[
 \det(D^{-T}MD^{-1})=\det M=\lambda^4\det M.
 \tag{OQC9}
\]
Thus $\lambda^4=1$ and $\lambda^5=1$, forcing $\lambda=1$.
If any nonidentity power fixed the line, its powers generate the
same cyclic group, so the preceding argument applies. This proves
the entire orbit classification.

There is an exact coefficient test, requiring no genericity premise.
Write the actual $Q=\sum_{|\mu|=2}q_\mu x^\mu$ and put
\[
 \mathfrak e(Q)=\sum_{\substack{|\mu|=2\\
                      \sum j\mu_j\not\equiv0\ (5)}}|q_\mu|^2.
 \tag{OQC10}
\]
All coefficients here are the entries of OQC3, including the full
inverse of the original period matrix. Since the monomials are a
basis and each summand is nonnegative, $\mathfrak e(Q)=0$ exactly
when $Q$ is invariant. The only degree-two weight-zero monomials
are $x_1x_4$ and $x_2x_3$. Hence the size-one case is exactly
$Q=a x_1x_4+b x_2x_3$ with $ab\ne0$; nonvanishing of the two
coefficients follows from OQC3. Positive $\mathfrak e(Q)$ gives
the five distinct irreducible orbit equations. These are exhaustive
outcomes of a test on the given periods, not assumptions on them.

In the size-one outcome, multiplication by $Q$ preserves every
weight. If $Qf$ is invariant, then $Q(gf-f)=0$ in the polynomial
ring, and $gf=f$. Therefore
\[
 \ker(\mathscr I_k\to R_k)=Q\mathscr I_{k-2},\qquad
 \operatorname{rank}\Lambda_k=d_{5,k,0}-d_{5,k-2,0},
 \quad \dim\ker\Lambda_k=q-d_{5,k,0}+d_{5,k-2,0}.
 \tag{OQC11}
\]
This also proves that $J_k=R_k^{\langle g\rangle}$: an invariant
class can be lifted to a polynomial, and averaging that lift gives
an invariant representative. The count itself follows from
\[
 d_{5,k,0}=\frac{\binom{k+3}{3}+4\epsilon_k}{5},\qquad
 \epsilon_k=\mathbf1_{k\equiv0\ (5)}-\mathbf1_{k\equiv1\ (5)}.
 \tag{OQC12}
\]
For a proof, the identity element has symmetric-power generating
series $(1-z)^{-4}$; each of the other four powers has series
$\prod_{j=1}^4(1-z\zeta^j)^{-1}=(1-z)/(1-z^5)$.
Averaging their coefficients gives OQC12. Subtraction and
$\binom{k+3}{3}-\binom{k+1}{3}=(k+1)^2$ evaluate the rank in
OQC11, for every $k\geq2$, as
\[
 \left\lfloor\frac{(k+1)^2}{5}\right\rfloor
                 +\mathbf1_{k\equiv0\text{ or }3\ (5)}.
 \tag{OQC13}
\]
The five numerator corrections before division by five are
$4,-4,-4,4,0$, respectively; the residues of $(k+1)^2$ are
$1,4,4,1,0$. This verifies every case in OQC13.

\subsection{An explicit invariant lift for the five-member orbit}
For the size-five outcome determined by OQC10, retain the full
actual products
\[
 \mathcal H_Q=\prod_{a=0}^4Q_a\quad(\deg\mathcal H_Q=10),
 \qquad \mathcal A_Q=\prod_{a=1}^4Q_a
                                      \quad(\deg\mathcal A_Q=8).
 \tag{OQC14}
\]
The action permutes all five factors cyclically, so
$g\mathcal H_Q=\mathcal H_Q$. None of the four factors of
$\mathcal A_Q$ is a scalar multiple of $Q$. Since each is
irreducible and $Q$ is prime, multiplication by $[\mathcal A_Q]$
is injective on the integral domain $R$.

For every homogeneous $f\in\mathscr P_d$ the exact trace identity is
\[
 \left[5\mathsf R(\mathcal A_Q f)\right]_{(Q)}
                          =[\mathcal A_Q f]_{(Q)}.
 \tag{OQC15}
\]
Indeed $g^a\mathcal A_Q$ is the product of all orbit factors
except $Q_a$. For $a\ne0$ it contains $Q_0=Q$. Those four
terms vanish in the displayed quotient, and the $a=0$ term is
exactly the right side. This proves OQC15 with its original
average and full factors, including every period-dependent
coefficient.

The lift also has a well-defined typed quotient map. Let
$\mathscr U=\mathbb C[x_1,x_2,x_3,x_4]/(\mathcal H_Q)$;
the invariant equation $\mathcal H_Q$ makes the cyclic action
well defined on this quotient. Then
\[
 \begin{gathered}
 \mathfrak L_d:R_d\longrightarrow
                      (\mathscr U_{d+8})^{\langle g\rangle},\qquad
 [f]_{(Q)}\longmapsto[5\mathsf R(\mathcal A_Q f)]_{(\mathcal H_Q)},\\
 \operatorname{res}\mathfrak L_d([f])=[\mathcal A_Q f]_{(Q)},
 \qquad \operatorname{res}:\mathscr U\to R.
 \end{gathered}\tag{OQC16}
\]
If $f$ changes by $Qb$, its proposed image changes by
$5\mathsf R(\mathcal H_Qb)=5\mathcal H_Q\mathsf Rb$, hence
by zero in $\mathscr U$. This proves well-definedness. The
image is invariant because it has an invariant polynomial
representative; OQC15 proves the second line. Injectivity follows
from the injectivity of multiplication by $[\mathcal A_Q]$ on
$R$. All maps in OQC16 are thus established, including their kernels.

The actual original invariant image consequently contains
\[
 [\mathcal A_Q]R_{k-8}\ \subseteq\ J_k\ \subseteq\ R_k.
 \tag{OQC17}
\]
The first inclusion follows by taking the invariant polynomial
representatives OQC15, and uses no dimension inference. Therefore,
for every $k\geq8$ in this outcome,
\[
 \operatorname{rank}\Lambda_k\geq(k-7)^2,\qquad
 \dim\ker\Lambda_k\leq(k+1)^2-(k-7)^2=16k-48.
 \tag{OQC18}
\]
These inequalities follow from the proved injection and OQC5.
They may be combined with the unconditional actual-matrix bounds
in OCS24, OCS31--33 by taking the stronger lower and upper bounds.
For the original degrees $k=9,13,17,21,25,29$, the displayed upper
bounds on the original kernel are respectively
$96,160,224,288,352,416$; at the first two degrees the earlier
unconditional bounds may be stronger. No such upper bound is
assigned to the size-one outcome, whose exact rank is OQC11.
Together OQC10--18 compute the complete orbit alternatives from
the original coefficients. Evaluating $\mathfrak e(Q)$ in the
original period family is a separate calculation on those explicit
coefficients; no outcome is postulated here.

\subsection{The induced original dual map and its complete action defect}
At $k\geq9$ the degree $d=k-8$ is positive, and the original
$d$-fold sum algebra is defined using the same $h,\Pi,U$.
For the size-five outcome, OQC7, OQC8 and OQC17 give the uniquely
typed injective map
\[
 \begin{gathered}
 \mathfrak J_k:E_{k-8}^\vee\hookrightarrow B_k^\vee,\\
 \Lambda_k^\vee\mathfrak J_k
        =\Theta_k\,M_{[\mathcal A_Q]}\,\Theta_{k-8}^{-1},\qquad
 \mathfrak J_k^\vee:B_k\twoheadrightarrow E_{k-8}.
 \end{gathered}\tag{OQC19}
\]
The right side of the second line has image in
$\operatorname{im}\Lambda_k^\vee$ by OQC17, and
$\Lambda_k^\vee$ is injective because $\Lambda_k$ is onto its
actual image. This gives existence and uniqueness. Each of the
three right-hand maps is injective, so $\mathfrak J_k$ is
injective. Finite-dimensional duality then proves the stated
surjection, using the canonical bidual identifications. The new
degree $k-8$ records contraction against the degree-eight
polynomial derived in OQC14. It does not replace either original
Gamma source of order $1$ or $k$.

The surviving action and its defect are fully determined. Put
\[
 \Omega=\operatorname{diag}(\gamma-i\delta,-\gamma-i\delta,
                              \gamma+i\delta,-\gamma+i\delta),
 \quad G=\mathsf H\Omega\mathsf H^{-1},\qquad
 (\mathcal Df)(x)=\sum_{j=1}^4(Gx)_j\partial_jf(x).
 \tag{OQC20}
\]
Here $x'(w)=Gx(w)$ by OQC1. In the original $t$ coordinates,
the weights of each product in $Q_0$ sum to zero, so
$\mathcal DQ=0$. Thus this derivation acts on each $R_d$.
Differentiating OQC7 proves exactly
\[
          \Theta_d\mathcal D=Y_d^\vee\Theta_d.
 \tag{OQC21}
\]
Both sides evaluated at $e^{wY_d}[1]$ are the derivative of
$f(x(w))$, and these vectors span $E_d$ by the distinct
exponentials; hence equality holds as maps, not only on a curve.
Leibniz's rule gives the complete defect
\[
 \begin{split}
 Y_k^\vee\Lambda_k^\vee\mathfrak J_k
       -\Lambda_k^\vee\mathfrak J_kY_{k-8}^\vee
 &=\Theta_k M_{[\mathcal D\mathcal A_Q]}\Theta_{k-8}^{-1},\\
 \mathcal D\mathcal A_Q
 &=\sum_{a=1}^4(\mathcal DQ_a)
                  \prod_{\substack{1\leq b\leq4\\b\ne a}}Q_b.
 \end{split}\tag{OQC22}
\]
The codomain in the first line is $E_k^\vee$. Its image has not
been asserted to lie in $\operatorname{im}\Lambda_k^\vee$.
Restriction to $\ker\Lambda_k$ supplies the exact remaining
map into $(\ker\Lambda_k)^\vee$, through OQC8. In particular
OQC19 is not used as an unproved intertwining of boundary actions.
OQC22 retains every term that can leave the observed dual subspace.

Finally, the failure of the cyclic average to commute with the
original action is explicitly calculable before any quotient:
\[
 ([\mathcal D,g^a]f)(x)
   =(\nabla f)(D^{-a}x)^T
                   (D^{-a}G-GD^{-a})x,
 \qquad
 [\mathcal D,\mathsf R]=\frac15\sum_{a=0}^4[\mathcal D,g^a].
 \tag{OQC23}
\]
The chain rule applied separately to $f(D^{-a}x)$ and to
$(\mathcal Df)(D^{-a}x)$ proves both identities. Thus all
action comparisons in the invariant lift have explicit original
coefficient maps. The rank estimates in OQC18 concern exactly
the kernel used in the original Gamma and arithmetic quotient
Grams, because those constructions use the same $\Lambda_k$.
An estimate for its volume still requires those original metrics;
no volume conclusion has been inferred solely from its dimension.


% Exact completion of the orbit rank calculation. No period outcome is assumed.
\section{Exact original rank from the complete quadric orbit}
\label{sec:OQX}

Retain every original object and map in OQC1--23. In particular
$m=1$, $k=4l+1$, $l\geq1$, $q=(k+1)^2$,
$\mathscr P=\mathbb C[x_1,x_2,x_3,x_4]$, and
$g f(x)=f(D^{-1}x)$ for the original Fourier matrix
$D=\operatorname{diag}(\zeta,\zeta^2,\zeta^3,\zeta^4)$.
The polynomial $Q=Q_0(\mathsf H^{-1}x)$ has its actual original
period and amplitude coefficients. Write
$\mathscr I=\mathscr P^{\langle g\rangle}$,
$R=\mathscr P/(Q)$ and $J=\operatorname{im}(\mathscr I\to R)$,
with the grading and the actual observation identification of OQC8.
For negative indices put $\mathscr I_d=0$ and $d_{5,d,0}=0$.

The orbit decision is the exhaustive test OQC10. The size-one
case already has the exact rank OQC11--13. The following calculation
completes the other case using the same original coefficients. It
does not assert that a period parameter belongs to that case before
its coefficient test has been evaluated.

\subsection{Every invariant vanishing polynomial contains the full orbit}

In the size-five case, let $Q_a=g^aQ$ and retain exactly
$\mathcal H_Q=\prod_{a=0}^4Q_a$ from OQC14. Each $Q_a$ is
irreducible, because an invertible linear substitution preserves
factorizations and $Q$ is irreducible by OQC3. The five factors are
pairwise nonassociate by the proved orbit classification, and each
is prime in the original polynomial ring. Consequently
\[
 \boxed{\quad
 \mathscr I\cap Q\mathscr P=\mathcal H_Q\mathscr I,
 \qquad
 \ker(\mathscr I_d\longrightarrow R_d)
                       =\mathcal H_Q\mathscr I_{d-10}.
 \quad}\tag{OQX1}
\]
Here is the complete divisibility proof. If $f\in\mathscr I$ and
$f=Qb$, then for every $a=0,\ldots,4$ the exact equality
$g^af=f$ gives $f=Q_a g^ab$. Thus every one of the five original
prime factors divides $f$. Distinctness of the factors implies
that their product divides $f$: after factoring out $Q_0$,
primality of $Q_1$ and the fact that $Q_1$ does not divide $Q_0$
force $Q_1$ to divide the remaining quotient. Repeating this
argument with $Q_2,Q_3,Q_4$ proves
$f=\mathcal H_Q b_0$. The action permutes the five exact factors
cyclically, so $g\mathcal H_Q=\mathcal H_Q$. Applying $g$ to
this last equality and using $gf=f$ gives
$\mathcal H_Q(gb_0-b_0)=0$. The polynomial ring is an integral
domain and $\mathcal H_Q\ne0$, whence $gb_0=b_0$.
This proves the first inclusion in OQX1. Conversely an invariant
$b_0$ gives an invariant $\mathcal H_Qb_0$ divisible by $Q$,
which proves the reverse inclusion. The product is homogeneous
of degree ten. If its product with $b_0$ is homogeneous of degree
$d$, decomposition of $b_0$ into homogeneous parts shows that
only its part of degree $d-10$ can be nonzero: each other part
has zero product with $\mathcal H_Q$ and hence is zero.
This proves the graded kernel formula, including $d<10$.

\subsection{The exact invariant quotient and the restriction map}

Let $\mathscr U=\mathscr P/(\mathcal H_Q)$ as in OQC16. The
following are isomorphisms of graded complex algebras with their
specified maps:
\[
 \boxed{\quad
 \mathscr I/(\mathcal H_Q\mathscr I)
 \xrightarrow{\ [f]\mapsto[f]_{(\mathcal H_Q)}\ }
       \mathscr U^{\langle g\rangle}
 \xrightarrow{\ \operatorname{res}\ } J\subseteq R.
 \quad}\tag{OQX2}
\]
The action on $\mathscr U$ is defined because $\mathcal H_Q$ is
invariant. If $[f]_{(\mathcal H_Q)}$ is invariant, averaging its
representatives gives
$[\mathsf R f]_{(\mathcal H_Q)}=[f]_{(\mathcal H_Q)}$, since
$[g^af]=[f]$ for every $a$. Thus the first map is onto.
Its kernel before passage to the displayed domain is
$\mathscr I\cap\mathcal H_Q\mathscr P
=\mathcal H_Q\mathscr I$: the reverse inclusion is immediate,
and for $f=\mathcal H_Qb\in\mathscr I$, cancellation in
$\mathcal H_Q(gb-b)=0$ proves $b\in\mathscr I$.
This proves that the first map is injective as well.

The full quotient restriction
$\operatorname{res}:\mathscr U\to R$ is well defined because
$(\mathcal H_Q)\subseteq(Q)$; its kernel is exactly
$Q\mathscr P/\mathcal H_Q\mathscr P$.
For its invariant subspace, choose the invariant representative
just constructed. If its restriction vanishes, that representative
belongs to $\mathscr I\cap Q\mathscr P$, which is
$\mathcal H_Q\mathscr I$ by OQX1. Its class in $\mathscr U$ is
therefore zero. The invariant restriction is injective. Its image
is exactly $J$, since an invariant representative restricts to
$J$ by definition and every element of $J$ has such a representative.
Both maps are induced by polynomial quotient maps, so they preserve
products, sums, complex scalars and degrees. No action of $g$ on
$R$ is asserted in the size-five case. In particular the second
map retains precisely the invariant image of the original
observation, with its complete restriction kernel already computed.

\subsection{All original degrees and the exact kernel dimension}

The map $\Theta_k$ and the original observation in OQC8 identify
$\dim J_k$ with $\operatorname{rank}\Lambda_k$ and identify
$R_k/J_k$ with $(\ker\Lambda_k)^\vee$. Multiplication by the
nonzero polynomial $\mathcal H_Q$ is injective in $\mathscr P$.
Taking dimensions in OQX1 thus gives, for every nonnegative degree
$d$, the exact invariant restriction rank
$\dim J_d=d_{5,d,0}-d_{5,d-10,0}$. At the original degrees this
evaluates to
\[
 \boxed{\quad
 \operatorname{rank}\Lambda_k
      =d_{5,k,0}-d_{5,k-10,0}
      =k^2-6k+17,\qquad
 \dim\ker\Lambda_k=8k-16,
 \quad k=4l+1,\ l\geq1,
 \quad}\tag{OQX3}
\]
in the size-five outcome of the actual coefficient test.
To verify every constant and degree range, for $k\geq10$ the
complete character formula OQC12 gives
\[
 d_{5,k,0}-d_{5,k-10,0}
 =\frac15\left\{\binom{k+3}{3}-\binom{k-7}{3}
                         +4(\epsilon_k-\epsilon_{k-10})\right\}.
\]
The two $\epsilon$ values are equal because their indices differ
by ten. Expanding the remaining cubics gives
\[
 \binom{k+3}{3}-\binom{k-7}{3}
 =\frac{(k^3+6k^2+11k+6)
             -(k^3-24k^2+191k-504)}6
 =5(k^2-6k+17).
\]
The only original degrees below ten are five and nine. Their
negative-degree terms in OQX1 vanish, while the same exact
character formula gives
\[
 d_{5,5,0}=(56+4)/5=12=5^2-6\cdot5+17,
 \qquad
 d_{5,9,0}=220/5=44=9^2-6\cdot9+17.
\]
This proves the first equality of OQX3 for every original $l\geq1$.
Since $\dim E_k=\dim R_k=(k+1)^2$, rank-nullity gives
$(k+1)^2-(k^2-6k+17)=8k-16$, proving the second equality.
For the original degrees already tabulated in OCS26, the exact
five-orbit values are
\[
 \begin{array}{c|r|r}
 k&\operatorname{rank}\Lambda_k&\dim\ker\Lambda_k\\\hline
 5&12&24\\
 9&44&56\\
 13&108&88\\
 17&204&120\\
 21&332&152\\
 25&492&184\\
 29&684&216
 \end{array}
\]
These are exact values for this exhaustively classified orbit
outcome, replacing the less precise bounds OQC18 wherever that
outcome occurs. The full invariant lift and its action defect
OQC15--23 remain valid with all original coefficients. The
size-one outcome retains its different exact values OQC11--13.
No choice between the two outcomes has been made here without
evaluating the original period coefficient test, and no estimate
for a metric volume has been inferred from its dimension alone.

\end{document}
